\documentclass{article}
\usepackage{PRIMEarxiv} % Core package for the PrimeAI Template
\usepackage{amsmath, amssymb, amsthm, bm, dcolumn} % Add amsthm here for the proof environment
\usepackage[numbers,sort&compress]{natbib} % Natbib for citations
\usepackage{graphicx} % For high-quality images
\usepackage{multicol} % Multi-column figures/tables
\usepackage{hyperref} % Hyperlinks
\usepackage{listings} % Code listings
\usepackage{authblk} % For structured affiliations
% Define theorem style (optional)
\theoremstyle{plain}
\newtheorem{theorem}{Theorem}
\renewcommand{\qedsymbol}{}

% Custom commands for primes and qubit encoding
\newcommand{\primeQubit}[1]{|p_{#1}\rangle}
\newcommand{\primeGate}[1]{U_{p_{#1}}}

\usepackage{wrapfig}
\usepackage[pscoord]{eso-pic}
\usepackage[fulladjust]{marginnote}
\reversemarginpar

% Typesetting improvements without footnote patching
\usepackage[protrusion=true, expansion=true, tracking=false]{microtype}
\microtypecontext{spacing=nonfrench}

% Line numbers
\usepackage[right]{lineno}

% Text layout - adjust as needed
\raggedright
\setlength{\parindent}{0.5cm}
\textwidth 5.25in 
\textheight 8.75in

% Set double spacing
\usepackage{setspace} 
\doublespacing

% Adjust width for specific content
\usepackage{changepage}

% Adjust caption style
\usepackage[aboveskip=1pt,labelfont=bf,labelsep=period,singlelinecheck=off]{caption}

% Remove brackets from references
\makeatletter
\renewcommand{\@biblabel}[1]{\quad#1.}
\makeatother

% Header, footer, and page numbers
\usepackage{lastpage,fancyhdr}
\pagestyle{fancy}
\fancyhf{}
\fancyhead[L]{Preprint - PrimeAI Enhanced Template}
\fancyfoot[C]{\scriptsize Multiplicity Theory © 2025 Ryan O. Van Gelder - Citizen Gardens \\ Licensed Under MIT and CC BY-NC-SA 4.0.}
\fancyfoot[R]{Page \thepage\ of \pageref{LastPage}}
\renewcommand{\footrule}{\hrule height 2pt \vspace{2mm}}

\begin{document}
\title{Enhancing Coupled Harmonic Oscillators with Multiplicity Theory and Hyperelliptic Curves}

\author{Ryan O. Van Gelder}
\affil{Citizen Gardens - The Foundation of Multiplicity}
\date{\today}
\maketitle


\begin{abstract}
This work explores the application of Multiplicity Theory to improve the topology of coupled harmonic oscillators influenced by periodic electric and perpendicular magnetic fields. By integrating hyperelliptic curve frameworks, we provide a comprehensive mathematical model to analyze energy flows, stability, and dynamic coupling. This approach leverages recursive feedback, tensor networks, and prime-based encoding, extending applications to quantum computing, signal processing, and advanced materials.
\end{abstract}

\section{Introduction}
Coupled harmonic oscillators under external fields provide an essential model for studying energy transfer, resonance, and stability in physical systems. Multiplicity Theory offers a prime-based encoding and recursive feedback mechanism to model nonlinear dynamics effectively. The inclusion of hyperelliptic curves further enhances the geometric representation of these systems, enabling deeper insights into their topological and dynamic properties.

\section{Theoretical Foundation}

\subsection{Multiplicity Theory}
Multiplicity Theory integrates eigenvalue dynamics, tensor coupling, and phase coherence to describe interconnected systems. The dynamic multiplicity equation incorporates time-evolving parameters:\newline
\begin{equation}
\frac{\partial \rho_k}{\partial t} = \alpha_k(t) \rho_k + \beta_k(t) I_k + \gamma_k(t) \sum_j T_{kj} \rho_j + \lambda(t)\big(\Omega_B(\rho) + \Omega_{FS}(\rho)\big) + \eta_k \rho_k^2 + \xi_k(t),
\end{equation}
where $\rho_k$ represents the state variable, $\Omega_B$ and $\Omega_{FS}$ are geometric feedback terms, and $\xi_k(t)$ captures stochastic noise.

\subsection{Coupled Harmonic Oscillators}
Harmonic oscillators interacting with periodic electric fields $E(t) = E_0 \sin(\omega t)$ and perpendicular magnetic fields $B(t)$ demonstrate energy transfer and phase shifts:\newline
\begin{equation}
H(t) \ni \psi(t) \rightarrow M(t,\psi(t)) T(t,\psi(t)) + f(t,\psi(t)) = \lambda(t) \psi(t).
\end{equation}
The tensor $T(t,\psi(t))$ models coupling dynamics, while $f(t,\psi(t))$ introduces external driving forces.

\subsection{Hyperelliptic Curves}
Hyperelliptic curves, described as branched double covers, capture topological branching induced by magnetic effects:\newline
\begin{equation}
y^2 = f(x), \quad \text{where } f(x) \text{ defines the potential landscape.}
\end{equation}
Branch points represent repeated roots, aligning with degenerate states in harmonic oscillators.

\section{Applications of Multiplicity Theory}

\subsection{Energy and Stability Dynamics}
Multiplicity Theory models energy transfer and stability in oscillators by incorporating:\newline
\begin{itemize}
    \item Tensor networks to represent multi-dimensional interactions.
    \item Prime-based encoding for efficient state mapping and coupling.
    \item Recursive feedback for adaptive state evolution.
\end{itemize}

\subsection{Time-Evolution Equations}
The system's evolution is governed by:
\begin{equation}
H(t) \ni \psi(t) \rightarrow M(t,\psi(t)) T(t,\psi(t)) + f(t,\psi(t)) = \lambda(t) \psi(t).
\end{equation}
Dynamic parameters ensure real-time adaptability.

\subsection{Branching and Covering}
Magnetic field-induced branching in hyperelliptic curves is modeled using recursive frameworks:\newline
\begin{equation}
\lambda(t) = \sum_{i=1}^N \lambda_i \mu_i e^{i\theta_i(t)} v_i.
\end{equation}
This approach elucidates topological bifurcations and their impact on energy landscapes.

\subsection{Potential Landscape Analysis}
Potential functions with repeated roots introduce resonance and degeneracies. Multiplicity Theory models these phenomena using tensor representations:
\begin{equation}
V(x) = \sum_{i=1}^n a_i (x - x_i)^2.
\end{equation}

\section{Numerical Simulations and Applications}

\subsection{Quantum Computing and Signal Processing}
Tensor networks and prime-based encoding optimize quantum algorithms for state evolution and signal processing. For a given quantum state $\psi(t)$, represented as:
\begin{equation}
\psi(t) = \sum_{i=1}^N c_i(t) \phi_i,
\end{equation}
the coefficients $c_i(t)$ are dynamically updated using recursive feedback equations:\newline
\begin{equation}
\frac{dc_i}{dt} = \alpha_i c_i + \sum_{j=1}^N T_{ij} c_j + \beta_i \sin(\omega t).
\end{equation}
This enables efficient simulation of state transitions under external fields.

\subsection{Advanced Materials}
Understanding energy flow and stability in coupled systems aids the design of materials with tailored properties. By modeling the potential landscape $V(x)$, the interaction energy can be computed as:
\begin{equation}
U = \int \nabla V(x) \cdot dx,
\end{equation}
where $\nabla V(x)$ captures local stability and energy distribution. Tensor networks further refine these calculations to include multi-scale interactions.

\subsection{Simulation of Topological Effects}
Branching phenomena in hyperelliptic curves are simulated using eigenvalue dynamics:\newline
\begin{equation}
\lambda_i(t) = \lambda_i(0) e^{-\gamma_i t} + \int_0^t \kappa_i(t') dt',
\end{equation}
where $\gamma_i$ is the damping factor and $\kappa_i(t')$ represents external coupling effects. These simulations predict bifurcations and stability thresholds.


\section{Conclusion and Future Work}
This framework integrates Multiplicity Theory and hyperelliptic geometry to enhance the analysis of coupled harmonic oscillators. Future research will focus on experimental validation and extending applications to quantum systems and AI-driven simulations.


\nocite{*}
\bibliographystyle{unsrt}
\bibliography{references}

\end{document}